\section{Empirical Methodology}
\label{sec:methodology}

We evaluate \systemname across seven benchmark arenas spanning software repair,
GPU-kernel optimization, language-model training, training-speed optimization,
research-assistant tasks, and data synthesis. Each benchmark is reported in its
native unit. SWE-Bench Pro is one benchmark in this suite; its sequential task log
also enables deeper analyses of runtime evolution and Reviewer intervention.

\subsection{Research Questions}

Evaluation is organized around five research questions, labelled RQ to keep them
distinct from the Driver questions Q1--Q4 of Section~\ref{sec:introduction}, which
are about what the runtime decides rather than what we measure. RQ0 concerns the
operating mode itself and is prerequisite to the rest: a result obtained under
continuous supervision and the same result obtained unattended are different
findings.

\paragraph{RQ0: Unattended operation.}
How often does the runtime require a human, what does it ask for when it does, and
what fraction of available time is it producing? These are reported as duty cycle and
interruption rate over campaign ledgers, together with what the unattended hours
delivered.

\paragraph{RQ1: Benchmark performance.}
What task-native outcomes does \systemname achieve across the seven benchmark
arenas, and how do they compare with the reference reported for each arena?

\paragraph{RQ2: Fixed-model runtime self-evolution.}
Within the SWE-Bench Pro run, how does solve-time resource demand change as reviewed
Skill, Wiki, verification, and routing state accumulates while the underlying model
parameters remain fixed?

\paragraph{RQ3: Reviewer routing and recovery.}
Within the same SWE-Bench Pro run, how often is an independent Reviewer invoked,
and how many initially rejected solutions are recovered after Reviewer feedback?

\paragraph{RQ4: Process-to-capability mechanisms.}
How do the retained traces quantify process-data advantage, verification-gated
correction, Token conversion, and cross-task reuse?

\subsection{Measuring Unattended Operation}

Duty cycle and interruption rate are computed from two runtime ledgers rather than
from instrumentation added for this report. The settled usage ledger records one row
per model call with start and completion timestamps, so the union of call intervals
gives occupied time without relying on a silence heuristic; a single call can exceed
two hours, which is why event-gap thresholds are unsuitable. The event log records
system-initiated human requests as typed events, which separates them from operator
messages that the system never waited on. Campaigns with fewer than 50 recorded calls
are excluded as too short to characterise. The two subtracted terms in
Equation~\ref{eq:duty-cycle}---empty backlog and unusable environment---are
identified from gaps exceeding one hour, with the distinction drawn by whether the
preceding event was a completed mission or an interrupted execution. Because these
are production campaigns rather than a controlled protocol, task difficulty, vertical,
and operator availability all vary across rows; the tables characterise an operating
profile and are not a matched comparison.

\subsection{Benchmark Suite}

SWE-Bench Pro evaluates repository-level software-engineering tasks with executable
acceptance tests~\citep{deng2025swebenchpro}. The remaining arenas evaluate B200
kernel optimization, nanochat training on B200 and H100, the nanoGPT speedrun,
AARRI-Bench, and Math-Reasoning Data Synthesis from the Arbor suite
~\citep{lin2026solexecbench,karpathy_nanochat,karpathy_autoresearch,
jordan_nanogptbench,wang2026aarri,jin2026arbor}. Because rank, bits per byte,
elapsed time, solve rate, and pass-gap are incompatible metrics, we do not average
them into one score.

\begin{table}[H]
\centering
\footnotesize
\renewcommand{\arraystretch}{1.16}
\begin{tabularx}{\linewidth}{@{}p{2.15cm} X p{3.25cm} p{1.05cm}@{}}
\toprule
\rowcolor{papersand}
\textbf{Benchmark} & \textbf{What the agent must do} & \textbf{Evaluation metric} & \textbf{Better} \\
\midrule
SWE-Bench Pro & Repair a real repository issue by editing production code; the submitted patch must satisfy task-specific executable acceptance tests. & Fraction of tasks resolved by the official verifier & Higher \\
SOL-ExecBench & Optimize real GPU kernels without changing their outputs, targeting execution close to an analytically derived hardware speed-of-light bound. & SOL score, global rank, and high placements after correctness checks & Higher \\
nanochat B200 & Improve the nanochat training program under a fixed five-minute run on one B200 GPU. & Validation bits per byte (BPB), a compression-oriented language-model metric & Lower \\
nanochat H100 & Solve the same five-minute nanochat training objective on one H100; it is a separate hardware-specific optimization problem. & Validation BPB under the H100 protocol & Lower \\
nanoGPT speedrun & Modify the training system so an eight-H100 run reaches validation loss $3.28$ as quickly as possible. & Mean time to the target over ten verifier runs & Lower \\
AARRI-Bench & Complete 82 granular tasks representing the judgment, diligence, and professional practice expected from a research intern. & Number and percentage of tasks solved & Higher \\
Math-Reasoning Data & Build a pipeline that generates AIME-style reasoning problems that are difficult on the first attempt but solvable with additional attempts. & Mean $\mathrm{pass@4}-\mathrm{pass@1}$ gap & Higher \\
\bottomrule
\end{tabularx}
\caption{Definitions of the seven benchmark arenas. B200 and H100 nanochat runs
are kept separate because hardware changes the optimization problem.}
\label{tab:benchmark-definitions}
\end{table}

The six non-SWE arenas use GPT-5.5 through Codex. The SWE-Bench Pro evaluation uses
GPT-5.5/xhigh through Copilot for both Direct Copilot and \systemname. The full
731-task comparison reports approximate accuracy and the aggregate Token ratio
\begin{equation}
  R_{\mathrm{tok}}
  = \frac{N_{\mathrm{\systemname}}^{\mathrm{total}}}
         {N_{\mathrm{Copilot}}^{\mathrm{total}}}
  \approx 1.41.
  \label{eq:aggregate-token-ratio}
\end{equation}
Raw Copilot Token totals and per-Wave resource traces were not retained; therefore
the longitudinal analysis is \systemname-only.

\subsection{SWE-Bench Pro Runtime Self-Evolution Analysis}

For a completed \systemname Wave $w$ with task set $\mathcal{T}_w$, we measure
\begin{align}
  \bar I_w^{\mathrm{solve}}
  &= \frac{1}{|\mathcal{T}_w|}
     \sum_{i\in\mathcal{T}_w} I_i^{\mathrm{solve}},
  \label{eq:solve-token-wave}\\
  \bar T_w^{\mathrm{agent}}
  &= \frac{1}{|\mathcal{T}_w|}
     \sum_{i\in\mathcal{T}_w} T_i^{\mathrm{agent}}.
  \label{eq:solve-time-wave}
\end{align}
$I_i^{\mathrm{solve}}$ contains all model calls in the accepted task trajectory,
including routing, Skill adaptation, execution, and any independent review.
$T_i^{\mathrm{agent}}$ measures active
workflow time. Orchestration wait, environment preparation, external verification,
infrastructure recovery, and post-task knowledge maintenance are excluded.

Completed Waves are summarized by task-weighted windows: W1--6 (startup), W7--12
(early reuse), W13--18 (composition shift), W19--22 (mature operation), and
W23--24 (late difficult tasks). Two incomplete Waves are omitted from the grouped
means. The primary comparison is startup versus mature operation; composition and
difficulty variation remain visible.

The unit of analysis is the observed sequential window, not a matched task pair.
The startup window begins with less project-specific state, whereas later windows
can retrieve reviewed repository knowledge, reusable procedures, prior failure
records, verification guidance, and updated routing decisions. Lower Token or time
demand in a later window is therefore consistent with runtime self-evolution when
the retained state removes repeated inspection or failed work. It is not, by
itself, a causal estimate of the value of that state. Task identity, repository
mix, execution latency, and difficulty also change across Waves, and no matched
frozen-state replay is available.

\subsection{SWE-Bench Pro Reviewer Analysis}

A task is \emph{Reviewer-invoked} when its final trajectory contains an independent
Reviewer model call; otherwise it follows Engineer self-review. A Reviewer verdict
of \code{continue} is a revision request. We report both subsequent official
verifier success and the stricter case in which a later Reviewer verdict is
\code{done}. The latter is called a \emph{strict Reviewer rescue}.

Reviewer routing is adaptive rather than randomized. The analysis reports routing,
revision, and recovery behavior over the observed task sequence.

The 41-artifact research portfolio is reported separately as a measure of research
program breadth.

\subsection{Representative Vertical-Trace Protocol}

To complement endpoint and software-trajectory measurements, we report one
Erd\H{o}s--Gy\'arf\'as mathematical campaign as a representative vertical trace.
The trace is role-led rather than theorem-led: it records what the Manager,
Planner, Engineer, and Reviewer did during recovery, problem selection, route
testing, proof production, revision, acceptance, and state retention. The reported
record contains one accepted route falsification and six proof-backed research
deltas; detailed mathematical statements remain in the supplementary claim table.

The protocol tracks whether the runtime retains a falsified branch, carries named
blockers across missions, and admits later work after artifact-based review. The
role trace, claim table, and figure provenance are released with the
supplementary materials.

\subsection{Theory-to-Measurement Mapping}

We map the theory to four reported quantities. Immediate information gain
$\Delta I_t$ is represented only in task-native units: benchmark improvement,
accepted theorem delta, or certified branch elimination. Review correction is
represented by Reviewer \code{continue}--revise trajectories and later
external-verifier outcomes. Reuse value $G_L$ is approximated observationally by
later Token/time demand as Skill and Wiki state accumulates. Token conversion is
represented by the role- and mission-level attribution in the mathematical case.
Appendix~\ref{app:formula-instantiation} records the corresponding substitutions in
their native units.
